% ============================================================
% CAPITOLO 5 - LA SCRITTURA A QUATTRO PARTI
% ============================================================
\chapter{La scrittura a quattro parti}
\label{chap:condotta_voci}

Gli accordi studiati finora sono entità verticali. Per usarli bisogna distribuire le loro note fra quattro voci e collegarli in modo fluido. Il modello di studio è la scrittura corale: Soprano, Contralto, Tenore e Basso. Questo capitolo tratta come si dispongono le voci; il collegamento fra accordi e gli errori da evitare sono nei Capitoli \ref{chap:collegamento} e \ref{chap:errori_condotta}.

\section{Disposizione delle voci}
\label{sec:disposizione_voci}

La scrittura a quattro parti simula un coro misto. A ogni voce si assegna un'estensione tipica, da rispettare per avere buona sonorità e cantabilità. Le note sono indicate con il numero d'ottava (Do4 è il Do centrale).

\begin{definizione}[Estensione delle voci corali (SATB)]
	\label{def:estensione_voci}
	Le estensioni standard e la zona preferibile di ciascuna voce sono:
	\begin{center}
	\begin{tabular}{lcc}
	\toprule
	\textbf{Voce} & \textbf{Estensione} & \textbf{Zona preferibile (medio-alta)} \\
	\midrule
	Soprano  & Do4--La5 & Sol4--Sol5 \\
	Contralto & Sol3--Re5 & Do4--Do5 \\
	Tenore   & La2--Fa$\sharp$4 & Mi3--Mi4 \\
	Basso    & Re2--Do4 & Fa2--La3 \\
	\bottomrule
	\end{tabular}
	\end{center}
\end{definizione}

Ogni voce canta più comoda e più chiara nella parte medio-alta della sua estensione. Gli estremi gravi sono poveri di suono e quelli acuti sono faticosi, quindi si usano solo quando la condotta delle parti lo richiede. Per questo gli esempi del manuale tengono le voci nella zona preferibile ed evitano gli estremi, salvo l'ultima nota di una linea che deve raggiungere un'altezza obbligata.

Oltre all'estensione conta la \textbf{disposizione} delle tre voci superiori (soprano, contralto, tenore). La loro distanza reciproca determina l'aspetto dell'accordo.

\begin{definizione}[Posizione stretta e posizione lata]
	\label{def:posizione_stretta_lata}
	\begin{itemize}
		\item \textbf{Posizione stretta:} la distanza tra soprano e tenore è inferiore a un'ottava; le tre voci superiori sono ravvicinate.
		\item \textbf{Posizione lata (o larga):} la distanza tra soprano e tenore è uguale o superiore a un'ottava.
	\end{itemize}
\end{definizione}

\begin{regola}[Distanza e separazione tra le voci]
	\label{reg:distanza_voci}
	Le voci devono restare unite ma distinguibili. Valgono questi limiti e queste preferenze:
	\begin{itemize}
		\item soprano-contralto e contralto-tenore: al massimo \textbf{un'ottava}, e di norma almeno una \textbf{terza}; una seconda tra voci adiacenti rende l'accordo confuso;
		\item tenore-basso: più libera, ma al massimo \textbf{una dodicesima}, e di norma almeno una \textbf{quinta}; un basso troppo vicino al tenore dà un suono pesante e poco chiaro;
		\item il basso sta nettamente sotto le voci superiori, che si dispongono in modo equilibrato tra loro, evitando che due voci si addensino mentre la terza resta isolata.
	\end{itemize}
\end{regola}

\section{Le posizioni delle triadi}
\label{sec:posizioni_triadi}

Oltre alla disposizione (stretta o lata), che riguarda la distanza tra le voci, esiste la \textbf{posizione} della triade.

\begin{definizione}[Posizione della triade]
	\label{def:posizione_triade}
	La posizione di una triade è data dalla nota dell'accordo che si trova al soprano, la voce più acuta:
	\begin{itemize}
		\item \textbf{posizione di ottava (o fondamentale):} al soprano c'è la fondamentale;
		\item \textbf{posizione di terza:} al soprano c'è la terza;
		\item \textbf{posizione di quinta:} al soprano c'è la quinta.
	\end{itemize}
\end{definizione}

Le posizioni non vanno confuse con i rivolti (Capitolo \ref{chap:rivolti}): i rivolti dipendono dalla nota al \textit{basso}, le posizioni da quella al \textit{soprano}.

\section{I raddoppi}
\label{sec:raddoppi}

Si scrive a quattro parti ma gli accordi di triade hanno tre note: una di esse va raddoppiata.

\begin{regola}[Principio generale dei raddoppi]
	\label{reg:raddoppi}
	\begin{itemize}
		\item Di norma si raddoppia la \textbf{fondamentale}, soprattutto negli stati fondamentali.
		\item \textbf{Mai} raddoppiare la \textbf{sensibile} (VII grado): le due sensibili dovrebbero salire entrambe alla tonica e darebbero ottave parallele.
		\item \textbf{Mai} raddoppiare la \textbf{settima} degli accordi di settima.
		\item In alcuni casi conviene raddoppiare la \textbf{terza} (purché non sia la sensibile): nel primo rivolto delle triadi minori e della triade diminuita, e nel \sigla{II}{7}{} incompleto che va al V, dove evita le quinte parallele (Regola \ref{reg:ii7_raddoppio}).
		\item Il III grado della scala maggiore (la mediante, nota caratteristica del modo) si raddoppia con parsimonia.
	\end{itemize}
\end{regola}

\section{Tipologie di moto}
\label{sec:tipologie_moto}

Quando due voci passano da un accordo al successivo disegnano due movimenti melodici, e la loro relazione definisce il tipo di moto.

\begin{definizione}[Tipi di moto tra le parti]
	\label{def:tipi_moto}
	Tra due voci qualsiasi esistono tre tipi di moto:
	\begin{itemize}
		\item \textbf{Moto retto (o diretto):} le due voci vanno nella stessa direzione. Se mantengono lo stesso intervallo si parla di \textit{moto parallelo}.
		\item \textbf{Moto contrario:} le due voci vanno in direzioni opposte.
		\item \textbf{Moto obliquo:} una voce sta ferma mentre l'altra si muove.
	\end{itemize}
\end{definizione}

Una buona scrittura preferisce il \textbf{moto contrario}, che garantisce indipendenza e chiarezza alle parti.
